\section{Transport du bloc modulaire au quotient d'incidence}
\label{sec:incidence-transport-new}
Le bloc central d'APP et l'incidence des faces et des sommets fournissent deux réalisations d'une même involution. Leur relation s'exprime au moyen d'un opérateur d'entrelacement rectangulaire, avant toute interprétation métrique. La congruence diagonale adjacente
\[
 k^2\equiv(k+1)^2\pmod9
 \quad\Longleftrightarrow\quad 2k+1\equiv0\pmod9
\]
sélectionne \(k=4\) dans la fenêtre déclarée. Le bloc correspondant est
\[
 B_c=\begin{pmatrix}7&2\\2&7\end{pmatrix}
 =7I_2+2R_c=9P_++5P_-,\qquad
 R_c=\begin{pmatrix}0&1\\1&0\end{pmatrix},\quad
 P_\pm=\frac{I_2\pm R_c}{2}.
\]
La différence entre les deux valeurs propres est de quatre. Le transport
incidenciel détermine ensuite comment cette différence se conserve
lorsque les multiplicités des secteurs changent.

Les faces du cube sont indexées par \((i,\epsilon)\), avec \(i\in\{1,2,3\}\) et \(\epsilon\in\{-1,1\}\) ; ses sommets sont indexés par \(\nu\in\{-1,1\}^3\). Nous définissons
\[
 M_{(i,\epsilon),\nu}=\mathbf1_{\{\nu_i=\epsilon\}}.
\]
Soient \(R_F\) l'opposition des faces et \(R_V\) l'inversion des sommets \(\nu\mapsto-\nu\). Dans chaque entrée, on vérifie
\[
 \mathbf1_{\{\nu_i=-\epsilon\}}
 =\mathbf1_{\{(-\nu)_i=\epsilon\}},\qquad R_FM=MR_V.
\]
La matrice \(M\) transporte donc l’involution complète.

\begin{theorem}[Descente incidentielle du bloc central]
Soient \(B_F=7I_6+2R_F\) et \(B_V=7I_8+2R_V\). Alors
\[
 B_FM=MB_V.
\]
Si \(u=(1,1,1,1,1,1)^{\mathsf T}\),
\(P_u=uu^{\mathsf T}/6\) et \(P_-=(I_6-R_F)/2\), le quotient
\[
 Q_\square=\RR^6/\ker(MM^{\mathsf T})
 \simeq\RR u\oplus E_-
\]
a pour dimension \(1+3\), et l’opérateur induit est
\[
 \overline B_F=9P_u+5P_-.
\]
\end{theorem}
\begin{proof}
La première identité résulte de la multiplication de \(R_FM=MR_V\) par deux et de l'ajout de \(7M\). Chaque face contient quatre sommets ; deux faces adjacentes en partagent deux et deux faces opposées n'en partagent aucun. De ce décompte, on obtient, entrée par entrée,
\[
 MM^{\mathsf T}=12P_u+4P_-.
\]
Le premier projecteur est de rang un et le second de rang trois ; leurs images
sont orthogonales. Les valeurs propres sont \(12,4,0\), avec les multiplicités
\(1,3,2\). L’opposition agit comme \(+1\) dans \(\RR u\) et comme
\(-1\) dans \(E_-\), et conserve le noyau de la matrice de Gram. Ainsi,
\(B_F\) descend au quotient et agit par neuf et cinq dans ces deux
facteurs directs. Son polynôme caractéristique est
\((\lambda-9)(\lambda-5)^3\).
\end{proof}

L’assignation temporelle à l’axe \(\RR u\) fixe la forme signée
\(J_L=P_--P_u\). Sur le quotient, \(I_Q=P_u+P_-\), de sorte que
\[
 J_L=\frac{7I_Q-\overline B_F}{2}.
\]
Cette égalité relie la forme de signature \((3,1)\), avec la
convention temporelle indiquée, et l’opérateur modulaire transporté. La
matrice de Gram initiale conserve sa positivité ; la signature est introduite
dans la réalisation ultérieure au moyen du signe assigné à chaque secteur.
Les dilatations neuf et cinq demeurent celles du bloc : l’identité
affine détermine la forme et ne transforme pas le bloc en une isométrie.

\subsection{Agrégation arithmétique, norme réticulaire et fonctionnelle de chemin}
La conservation de l'entier précédant la réduction devient visible lorsque l'on additionne les deux feuillets. Les matrices d'agrégation modulaire sont
\[
 M_\Pi=\begin{pmatrix}4&5&6\\5&4&6\\6&6&9\end{pmatrix},\qquad
 M_\Sigma=\begin{pmatrix}5&6&4\\6&4&5\\4&5&6\end{pmatrix}.
\]
Leurs sommes sont cinquante et un et quarante-cinq. Lors de la reconstruction des
neuf éléments de chaque bloc, l’entier et sa retenue vérifient
\[
 2025-810=(459-405)+9(174-45)=54+9\cdot129.
\]
L'agrégat résiduel et le quotient constituent deux coordonnées de cette égalité. Dans la construction réticulaire ultérieure, la valeur cinquante-quatre apparaît comme la norme de la chambre marquée \(v_\alpha=4\rho_1+\rho_2+\cdots+\rho_{12}\), pour des vecteurs orthogonaux \(\rho_j^2=2\). Dans la réalisation matérielle, elle apparaît comme l'évaluation de la fonctionnelle de chemin \(D_1\) du profil \((80,54,6)\). Chaque apparition conserve son application d'origine : agrégation avec retenue, forme réticulaire et évaluation orientée de chemin. La coïncidence de leurs valeurs acquiert un contenu structurel grâce à ces applications, dont les domaines et les preuves sont développés dans leurs sections respectives.
